\documentclass[a4paper,12pt,twoside]{memoir}

% Castellano
\usepackage[spanish,es-tabla]{babel}
\selectlanguage{spanish}
\usepackage[utf8]{inputenc}
\usepackage[T1]{fontenc}
\usepackage{lmodern} % Scalable font
\usepackage{microtype}
\usepackage{placeins}

\RequirePackage{booktabs}
\RequirePackage[table]{xcolor}
\RequirePackage{xtab}
\RequirePackage{multirow}

% Links
\PassOptionsToPackage{hyphens}{url}\usepackage[colorlinks]{hyperref}
\hypersetup{
	allcolors = {red}
}

% Ecuaciones
\usepackage{amsmath}

% Rutas de fichero / paquete
\newcommand{\ruta}[1]{{\sffamily #1}}

% Párrafos
\nonzeroparskip

% Huérfanas y viudas
\widowpenalty100000
\clubpenalty100000

% Imágenes

% Comando para insertar una imagen en un lugar concreto.
% Los parámetros son:
% 1 --> Ruta absoluta/relativa de la figura
% 2 --> Texto a pie de figura
% 3 --> Tamaño en tanto por uno relativo al ancho de página
\usepackage{graphicx}
\newcommand{\imagen}[3]{
	\begin{figure}[!h]
		\centering
		\includegraphics[width=#3\textwidth]{#1}
		\caption{#2}\label{fig:#1}
	\end{figure}
	\FloatBarrier
}

% Comando para insertar una imagen sin posición.
% Los parámetros son:
% 1 --> Ruta absoluta/relativa de la figura
% 2 --> Texto a pie de figura
% 3 --> Tamaño en tanto por uno relativo al ancho de página
\newcommand{\imagenflotante}[3]{
	\begin{figure}
		\centering
		\includegraphics[width=#3\textwidth]{#1}
		\caption{#2}\label{fig:#1}
	\end{figure}
}

% El comando \figura nos permite insertar figuras comodamente, y utilizando
% siempre el mismo formato. Los parametros son:
% 1 --> Porcentaje del ancho de página que ocupará la figura (de 0 a 1)
% 2 --> Fichero de la imagen
% 3 --> Texto a pie de imagen
% 4 --> Etiqueta (label) para referencias
% 5 --> Opciones que queramos pasarle al \includegraphics
% 6 --> Opciones de posicionamiento a pasarle a \begin{figure}
\newcommand{\figuraConPosicion}[6]{%
  \setlength{\anchoFloat}{#1\textwidth}%
  \addtolength{\anchoFloat}{-4\fboxsep}%
  \setlength{\anchoFigura}{\anchoFloat}%
  \begin{figure}[#6]
    \begin{center}%
      \Ovalbox{%
        \begin{minipage}{\anchoFloat}%
          \begin{center}%
            \includegraphics[width=\anchoFigura,#5]{#2}%
            \caption{#3}%
            \label{#4}%
          \end{center}%
        \end{minipage}
      }%
    \end{center}%
  \end{figure}%
}

%
% Comando para incluir imágenes en formato apaisado (sin marco).
\newcommand{\figuraApaisadaSinMarco}[5]{%
  \begin{figure}%
    \begin{center}%
    \includegraphics[angle=90,height=#1\textheight,#5]{#2}%
    \caption{#3}%
    \label{#4}%
    \end{center}%
  \end{figure}%
}
% Para las tablas
\newcommand{\otoprule}{\midrule [\heavyrulewidth]}
%
% Nuevo comando para tablas pequeñas (menos de una página).
\newcommand{\tablaSmall}[5]{%
 \begin{table}
  \begin{center}
   \rowcolors {2}{gray!35}{}
   \begin{tabular}{#2}
    \toprule
    #4
    \otoprule
    #5
    \bottomrule
   \end{tabular}
   \caption{#1}
   \label{tabla:#3}
  \end{center}
 \end{table}
}

%
% Nuevo comando para tablas pequeñas (menos de una página).
\newcommand{\tablaSmallSinColores}[5]{%
 \begin{table}[H]
  \begin{center}
   \begin{tabular}{#2}
    \toprule
    #4
    \otoprule
    #5
    \bottomrule
   \end{tabular}
   \caption{#1}
   \label{tabla:#3}
  \end{center}
 \end{table}
}

\newcommand{\tablaApaisadaSmall}[5]{%
\begin{landscape}
  \begin{table}
   \begin{center}
    \rowcolors {2}{gray!35}{}
    \begin{tabular}{#2}
     \toprule
     #4
     \otoprule
     #5
     \bottomrule
    \end{tabular}
    \caption{#1}
    \label{tabla:#3}
   \end{center}
  \end{table}
\end{landscape}
}

%
% Nuevo comando para tablas grandes con cabecera y filas alternas coloreadas en gris.
\newcommand{\tabla}[6]{%
  \begin{center}
    \tablefirsthead{
      \toprule
      #5
      \otoprule
    }
    \tablehead{
      \multicolumn{#3}{l}{\small\sl continúa desde la página anterior}\\
      \toprule
      #5
      \otoprule
    }
    \tabletail{
      \hline
      \multicolumn{#3}{r}{\small\sl continúa en la página siguiente}\\
    }
    \tablelasttail{
      \hline
    }
    \bottomcaption{#1}
    \rowcolors {2}{gray!35}{}
    \begin{xtabular}{#2}
      #6
      \bottomrule
    \end{xtabular}
    \label{tabla:#4}
  \end{center}
}

%
% Nuevo comando para tablas grandes con cabecera.
\newcommand{\tablaSinColores}[6]{%
  \begin{center}
    \tablefirsthead{
      \toprule
      #5
      \otoprule
    }
    \tablehead{
      \multicolumn{#3}{l}{\small\sl continúa desde la página anterior}\\
      \toprule
      #5
      \otoprule
    }
    \tabletail{
      \hline
      \multicolumn{#3}{r}{\small\sl continúa en la página siguiente}\\
    }
    \tablelasttail{
      \hline
    }
    \bottomcaption{#1}
    \begin{xtabular}{#2}
      #6
      \bottomrule
    \end{xtabular}
    \label{tabla:#4}
  \end{center}
}

%
% Nuevo comando para tablas grandes sin cabecera.
\newcommand{\tablaSinCabecera}[5]{%
  \begin{center}
    \tablefirsthead{
      \toprule
    }
    \tablehead{
      \multicolumn{#3}{l}{\small\sl continúa desde la página anterior}\\
      \hline
    }
    \tabletail{
      \hline
      \multicolumn{#3}{r}{\small\sl continúa en la página siguiente}\\
    }
    \tablelasttail{
      \hline
    }
    \bottomcaption{#1}
  \begin{xtabular}{#2}
    #5
   \bottomrule
  \end{xtabular}
  \label{tabla:#4}
  \end{center}
}



\definecolor{cgoLight}{HTML}{EEEEEE}
\definecolor{cgoExtralight}{HTML}{FFFFFF}

%
% Nuevo comando para tablas grandes sin cabecera.
\newcommand{\tablaSinCabeceraConBandas}[5]{%
  \begin{center}
    \tablefirsthead{
      \toprule
    }
    \tablehead{
      \multicolumn{#3}{l}{\small\sl continúa desde la página anterior}\\
      \hline
    }
    \tabletail{
      \hline
      \multicolumn{#3}{r}{\small\sl continúa en la página siguiente}\\
    }
    \tablelasttail{
      \hline
    }
    \bottomcaption{#1}
    \rowcolors[]{1}{cgoExtralight}{cgoLight}

  \begin{xtabular}{#2}
    #5
   \bottomrule
  \end{xtabular}
  \label{tabla:#4}
  \end{center}
}



\graphicspath{ {./img/} }

% Capítulos
\chapterstyle{bianchi}
\newcommand{\capitulo}[2]{
	\setcounter{chapter}{#1}
	\setcounter{section}{0}
	\setcounter{figure}{0}
	\setcounter{table}{0}
	\chapter*{#2}
	\addcontentsline{toc}{chapter}{#2}
	\markboth{#2}{#2}
}

% Apéndices
\renewcommand{\appendixname}{Apéndice}
\renewcommand*\cftappendixname{\appendixname}

\newcommand{\apendice}[1]{
	%\renewcommand{\thechapter}{A}
	\chapter{#1}
}

\renewcommand*\cftappendixname{\appendixname\ }

% Formato de portada
\makeatletter
\usepackage{xcolor}
\newcommand{\tutor}[1]{\def\@tutor{#1}}
\newcommand{\course}[1]{\def\@course{#1}}
\definecolor{cpardoBox}{HTML}{E6E6FF}
\def\maketitle{
  \null
  \thispagestyle{empty}
  % Cabecera ----------------
\noindent\includegraphics[width=\textwidth]{cabecera}\vspace{1cm}%
  \vfill
  % Título proyecto y escudo informática ----------------
  \colorbox{cpardoBox}{%
    \begin{minipage}{.8\textwidth}
      \vspace{.5cm}\Large
      \begin{center}
      \textbf{TFG del Grado en Ingeniería Informática}\vspace{.6cm}\\
      \textbf{\LARGE\@title{}}
      \end{center}
      \vspace{.2cm}
    \end{minipage}

  }%
  \hfill\begin{minipage}{.20\textwidth}
    \includegraphics[width=\textwidth]{escudoInfor}
  \end{minipage}
  \vfill
  % Datos de alumno, curso y tutores ------------------
  \begin{center}%
  {%
    \noindent\LARGE
    Presentado por \@author{}\\ 
    en Universidad de Burgos --- \@date{}\\
    Tutores: \@tutor{}\\
  }%
  \end{center}%
  \null
  \cleardoublepage
  }
\makeatother

\newcommand{\nombre}{Álvaro Manjón Vara} %%% cambio de comando

% Datos de portada
\title{NutriMenu}
\author{\nombre}
\tutor{Raúl Marticorena Sánchez y Antonio Jesús Canepa Oneto}
\date{\today}

\begin{document}

\maketitle


\newpage\null\thispagestyle{empty}\newpage


%%%%%%%%%%%%%%%%%%%%%%%%%%%%%%%%%%%%%%%%%%%%%%%%%%%%%%%%%%%%%%%%%%%%%%%%%%%%%%%%%%%%%%%%
\thispagestyle{empty}


\noindent\includegraphics[width=\textwidth]{cabecera}\vspace{1cm}

\noindent D. nombre tutor, profesor del departamento de nombre departamento, área de nombre área.

\noindent Expone:

\noindent Que el alumno D. \nombre, con DNI dni, ha realizado el Trabajo final de Grado en Ingeniería Informática titulado título de TFG. 

\noindent Y que dicho trabajo ha sido realizado por el alumno bajo la dirección del que suscribe, en virtud de lo cual se autoriza su presentación y defensa.

\begin{center} %\large
En Burgos, {\large \today}
\end{center}

\vfill\vfill\vfill

% Author and supervisor
\begin{minipage}{0.45\textwidth}
\begin{flushleft} %\large
Vº. Bº. del Tutor:\\[2cm]
D. nombre tutor
\end{flushleft}
\end{minipage}
\hfill
\begin{minipage}{0.45\textwidth}
\begin{flushleft} %\large
Vº. Bº. del co-tutor:\\[2cm]
D. nombre co-tutor
\end{flushleft}
\end{minipage}
\hfill

\vfill

% para casos con solo un tutor comentar lo anterior
% y descomentar lo siguiente
%Vº. Bº. del Tutor:\\[2cm]
%D. nombre tutor


\newpage\null\thispagestyle{empty}\newpage




\frontmatter

% Abstract en castellano
\renewcommand*\abstractname{Resumen}
\begin{abstract}
En este primer apartado se hace una \textbf{breve} presentación del tema que se aborda en el proyecto.
\end{abstract}

\renewcommand*\abstractname{Descriptores}
\begin{abstract}
Palabras separadas por comas que identifiquen el contenido del proyecto Ej: servidor web, buscador de vuelos, android \ldots
\end{abstract}

\clearpage

% Abstract en inglés
\renewcommand*\abstractname{Abstract}
\begin{abstract}
A \textbf{brief} presentation of the topic addressed in the project.
\end{abstract}

\renewcommand*\abstractname{Keywords}
\begin{abstract}
keywords separated by commas.
\end{abstract}

\clearpage

% Indices
\tableofcontents

\clearpage

\listoffigures

\clearpage

\listoftables
\clearpage

\mainmatter
\capitulo{1}{Introducción}

En la actualidad, cada vez se le da mayor importancia al hecho de llevar una alimentación saludable y equilibrada, puesto que es uno de los pilares clave para una vida longeva y libre de enfermedades. La Universidad de Burgos es consciente de esto, y es por ello por lo que forma parte de la Red Española de Universidades Promotoras de la Salud, y desde hace años se encuentra desarrollando el proyecto de Aula Campus Saludable, una sección que se encarga de promocionar hábitos saludables a todos los miembros de la Comunidad Universitaria.

Dentro de este área hay varias iniciativas y campañas que se han llevado a cabo, y uno de estos proyectos fue el del desarrollo de un conjunto de aplicaciones, cuyo objetivo era el de informar a los consumidores de los distintos centros de restauración de la Universidad de Burgos de los componentes nutricionales que contienen los platos que forman los menús disponibles. Este conjunto de aplicaciones también permitía la gestión de estos platos y menús, así como de los propios centros de restauración, sus empresas y sus usuarios.

La iniciativa para este nuevo proyecto es la de modernizar y mejorar este trabajo, para que su despliegue e implementación sean lo más sencillos posibles, así como intentar mejorar las aplicaciones webs lo máximo posible, manteniendo su propósito original.

\capitulo{2}{Objetivos del proyecto}

Los objetivos que se buscan con el desarrollo de este proyecto son los siguientes:

\section{Objetivos generales}
\begin{itemize}
	\item Implementar una infraestructura que sea agnóstica al hardware en el que se ejecute.
	\item Separar la infraestructura en servicios bien definidos y aislados.
	\item Desarrollar un sistema de orquestación del despliegue y hacer que este sea sencillo y escalable.
	\item Buscar una forma de poder acceder a nuevos datos nutricionales sin tener que tenerlos almacenados en local de forma obligatoria.
	\item Unificar las dos aplicaciones web en una, mejorar su interfaz y experiencia de usuario para que sean más sencillas de usar y hacer que sean funcionales en dispositivos móviles.
	\end{itemize}
	
\section{Objetivos técnicos}
\begin{itemize}
\item Contenerizar todos los servicios mediante Docker y Docker Compose.
\item Reestructurar la lógica del modelo de datos para separarla de la aplicación web, manteniendo el patrón Modelo - Vista - Controlador (MVC).
\item Desarrollar una API REST usando Spring Boot para que esta sea la vía por la que la aplicación interactúe con los datos.
\item Implementar un flujo de integración continua mediante el uso de GitHub Actions para comprobar que la infraestructura es funcional en todo momento.
\item Desarrollar una aplicación web con React que cumpla con todos los requisitos de las dos aplicaciones previas.
\item Hacer un diseño responsive con Bootstrap, que funcione en cualquier tamaño de pantalla.
\item Implementar la API externa de Nutritionix como fuente de composición nutricional de los alimentos.
\item Hacer uso de GitHub para la gestión del control de versiones de este proyecto, así como la administración de las tareas y los sprints.

\end{itemize}
\capitulo{3}{Conceptos teóricos}

En aquellos proyectos que necesiten para su comprensión y desarrollo de unos conceptos teóricos de una determinada materia o de un determinado dominio de conocimiento, debe existir un apartado que sintetice dichos conceptos.

Algunos conceptos teóricos de \LaTeX \footnote{Créditos a los proyectos de Álvaro López Cantero: Configurador de Presupuestos y Roberto Izquierdo Amo: PLQuiz}.

\section{Secciones}

Las secciones se incluyen con el comando section.

\subsection{Subsecciones}

Además de secciones tenemos subsecciones.

\subsubsection{Subsubsecciones}

Y subsecciones. 


\section{Referencias}

Las referencias se incluyen en el texto usando cite \cite{wiki:latex}. Para citar webs, artículos o libros \cite{koza92}.


\section{Imágenes}

Se pueden incluir imágenes con los comandos standard de \LaTeX, pero esta plantilla dispone de comandos propios como por ejemplo el siguiente:

\imagen{escudoInfor}{Autómata para una expresión vacía}{.5}



\section{Listas de items}

Existen tres posibilidades:

\begin{itemize}
	\item primer item.
	\item segundo item.
\end{itemize}

\begin{enumerate}
	\item primer item.
	\item segundo item.
\end{enumerate}

\begin{description}
	\item[Primer item] más información sobre el primer item.
	\item[Segundo item] más información sobre el segundo item.
\end{description}
	
\begin{itemize}
\item 
\end{itemize}

\section{Tablas}

Igualmente se pueden usar los comandos específicos de \LaTeX o bien usar alguno de los comandos de la plantilla.

\tablaSmall{Herramientas y tecnologías utilizadas en cada parte del proyecto}{l c c c c}{herramientasportipodeuso}
{ \multicolumn{1}{l}{Herramientas} & App AngularJS & API REST & BD & Memoria \\}{ 
HTML5 & X & & &\\
CSS3 & X & & &\\
BOOTSTRAP & X & & &\\
JavaScript & X & & &\\
AngularJS & X & & &\\
Bower & X & & &\\
PHP & & X & &\\
Karma + Jasmine & X & & &\\
Slim framework & & X & &\\
Idiorm & & X & &\\
Composer & & X & &\\
JSON & X & X & &\\
PhpStorm & X & X & &\\
MySQL & & & X &\\
PhpMyAdmin & & & X &\\
Git + BitBucket & X & X & X & X\\
Mik\TeX{} & & & & X\\
\TeX{}Maker & & & & X\\
Astah & & & & X\\
Balsamiq Mockups & X & & &\\
VersionOne & X & X & X & X\\
} 

\capitulo{4}{Técnicas y herramientas}

\begin{comment}
	Esta parte de la memoria tiene como objetivo presentar las técnicas metodológicas y las herramientas de desarrollo que se han utilizado para llevar a cabo el proyecto. Si se han estudiado diferentes alternativas de metodologías, herramientas, bibliotecas se puede hacer un resumen de los aspectos más destacados de cada alternativa, incluyendo comparativas entre las distintas opciones y una justificación de las elecciones realizadas. 
No se pretende que este apartado se convierta en un capítulo de un libro dedicado a cada una de las alternativas, sino comentar los aspectos más destacados de cada opción, con un repaso somero a los fundamentos esenciales y referencias bibliográficas para que el lector pueda ampliar su conocimiento sobre el tema.
\end{comment}

% TODO: Hacer introducción de esta sección %

\section{Patrón de diseño Modelo-Vista-Controlador (MVC)}

\imagen{MVC/diagrama-mvc}{Diagrama del patrón Modelo-Vista-Controlador}{0.9}

La función del patrón de diseño conocido como \textbf{Modelo-Vista-Controlador} es la de organizar de forma estructurada los componentes fundamentales de un determinado \textit{sistema de software}, estableciendo relaciones entre ellos.

Como metodología de trabajo, fue propuesta por Trygve Reenskaug en 1979, basándose su utilidad en tres ideas fundamentales \cite{campusmvp:mvc}:

\begin{itemize}
\item La organización y clasificación de la información que se aporta al sistema.
\item La unificación y gestión de la lógica del sistema.
\item La presentación de datos al usuario mediante una interfaz comprensible y asequible.
\end{itemize}

Los tres componentes fundamentales de este patrón tienen misiones diferenciadas sobre la forma en cómo tratan la información \cite{easyappcode:mvc}: 

\begin{itemize}
\item \textbf{Modelo}: Es el componente cuya misión es únicamente la de gestionar el contenido de una base de datos, incluida su manipulación y utilización. Es la parte que se va a encargar de la \textbf{gestión de los datos}.
\item \textbf{Vista}: Es el componente que sirve para acceder al contenido de la base de datos y presentar los datos al usuario. Va a ser \textbf{la parte con la que va a interactuar el usuario}.
\item \textbf{Controlador}: Este componente es el que pone en conexión el modelo y la vista, gestionando y procesando las instrucciones que se reciben para obtener un resultado. Su forma básica de trabajar consiste en acceder a la base de datos, manipularlos con una determinada finalidad y presentar los resultados de esa manipulación de los datos. Se va a encargar de gestionar \textbf{la lógica y la gestión de los datos}.
\end{itemize}
El motivo por el cual he decidido utilizar este patrón es porque permite trabajar de una forma precisa y eficaz, haciendo que la estructura del sistema sea clara, ordenada y separada, lo que facilita su comprensión y posterior mantenimiento.  

\section{Contenerización}

Se conoce como contenerización a la tecnología de virtualización que permite ejecutar \textbf{contenedores}, entornos completamente aislados del resto de la máquina que los está ejecutando, y que contienen todo lo necesario (bibliotecas, código, archivos de configuración, dependencias...) para que se pueda ejecutar una aplicación en cualquier entorno.\cite{microsoft:contenedores}

A diferencia de entornos clásicos de virtualización como las \textbf{máquinas virtuales}, en los que se emula el hardware por completo y se debe disponer de una instalación completa del sistema operativo para funcionar, en los contenedores se permite que múltiples instancias compartan un mismo sistema operativo, a pesar de estar usando espacios de ejecución distintos. Esto se consigue usando características de Linux como los \textit{espacios de nombres del kernel} o los \textit{cgroups} \cite{medium:kernelspace}, puesto que los contenedores están basados en Linux, pero esto no impide su ejecución en distintas plataformas y sistemas operativos, incluido Windows.
 
\subsection{Docker}

Docker es una plataforma de código abierto que permite el desarrollo, ejecución y distribución de aplicaciones en contenedores.

\subsubsection{Arquitectura de Docker}

\imagen{Docker/docker_architecture}{Arquitectura de Docker}{0.75}

La arquitectura de Docker es una arquitectura \textbf{cliente-servidor} \cite{docker:overview-architecture}. El cliente de Docker, que es con lo que interactúa el usuario, habla con el \textit{daemon} haciendo llamadas API REST, y este se encarga de gestionar las peticiones y los objetos de Docker, como son las imágenes, los contenedores, los volúmenes y las redes. Finalmente, para conseguir las imágenes, el \textit{daemon} habla con \textit{Docker Registry}, el repositorio que contiene todas las imágenes necesarias para el despliegue. Por defecto se usa \textbf{Docker Hub}, un repositorio público en el que cualquiera puede subir y descargar imágenes, pero se puede usar también un repositorio privado.

\subsubsection{Imágenes}

Las imágenes son plantillas que contienen instrucciones y parámetros para crear contenedores de Docker. Se puede crear un contenedor usando directamente una imagen del \textit{registry}, reutilizando imágenes ya existentes como base para crear nuestras propias imágenes, o incluso creando nuestras propias imágenes desde cero. Las imágenes se crean usando archivos \textit{Dockerfile}, añadiendo en él las distintas dependencias, instrucciones y parámetros, y cada uno de estos elementos es una capa dentro de la imagen. De esta manera, cada vez que se modifica el \textit{Dockerfile} y se vuelve a construir la imagen, realmente sólo se vuelven a construir las capas en las que han habido cambios, acelerando de esta forma el proceso de despliegue.
\cite{docker:overview-objects}

\subsection{Docker Compose}

Docker Compose es una herramienta que nos permite definir, orquestrar, desplegar y escalar \textbf{aplicaciones Docker multi-contenedor}. \cite{docker-compose:overview} Para ello, se define un archivo \textit{YAML} en el que se van a definir las  configuraciones necesarias para los distintos servicios pertenecientes a la aplicación. Esto nos va a permitir establecer relaciones y dependencias entre distintos contenedores que forman parte de un mismo aplicativo, permitiendo así que varios contenedores formen parte de una misma red o tengan la capacidad de compartir volúmenes de almacenamiento, por poner unos ejemplos.

He decidido usar esta herramienta ya que facilita bastante la conexión e integración entre los distintos servicios que componen la aplicación, como son la base de datos, el backend, y el frontend. Además de esto, a la hora de poner la aplicación en producción, obtenemos varias ventajas respecto a un entorno tradicional:

\begin{itemize}

\item \textbf{Automatización de despliegues}, ya que una vez hayamos definido en el archivo de configuración de Docker Compose todas las dependencias, configuraciones y relaciones necesarias, con un sólo comando Docker Compose se va a encargar de crear y desplegar los contenedores, configurar los elementos necesarios, e iniciar los servicios.
\item \textbf{Creación de entornos aislados y consistentes}, lo que nos permite desplegar entornos de desarrollo locales, así como entornos CI/CD de validación y pruebas, y que estos repliquen la infraestructura de producción, que también puede ser desplegada usando Docker Compose.
\item \textbf{Escalabilidad de infraestructura}, puesto que con un sólo comando podemos aumentar o disminuir el número de instancias de las que dispone un servicio.	
\end{itemize}

\section{Frameworks}

% TODO: Hacer esta parte %

\subsection{Spring Boot}

% TODO: Hacer esta parte %

\subsection{React}

% TODO: Hacer esta parte %

\subsection{JPA}

% TODO: Hacer esta parte %

\subsection{Bootstrap}

% TODO: Hacer esta parte %

\section{Gestores de paquetes}

% TODO: Hacer esta parte %

\subsection{Maven}

% TODO: Hacer esta parte %

\subsection{Node}

% TODO: Hacer esta parte %

\section{API}

% TODO: Hacer esta parte %

\subsection{Nutritionix}

% TODO: Hacer esta parte %

\section{Lenguajes de programación}

% TODO: Hacer esta parte %

\subsection{HTML}

El lenguaje HTML (\textit{HyperText Markup Language}) hoy en día es la base de internet, ya que es prácticamente imposible encontrar una página web que no tenga definido su contenido usando este lenguaje. Como \textbf{lenguaje de marcado} que es, utiliza marcas (etiquetas) que indican qué tipo de contenido se está representando en el navegador, lo que permite identificar y estructurar mejor este contenido, así como ser más precisos posteriormente a la hora de dar estilo a este contenido \cite{mdn:html}.

Otro de los valores fundamentales del lenguaje HTML es el \textbf{hipertexto}, que es lo que permite enlazar fácilmente contenidos, ya sea dentro de la misma página web o con otras distintas. Esto permite mejorar la forma de acceso a la información e interconectar el conocimiento, algo muy importante puesto que este es uno de los pilares básicos de internet.

Un archivo HTML sólo contiene (o debería contener, según las buenas prácticas \cite{midudev:html}) \textbf{información referente a la estructura y el contenido de la página web}, pero no información relacionada con el cómo se va a ver estructurada esta información o cómo se van a procesar los datos y se van a interactuar con ellos, puesto que eso es tarea de otros lenguajes, como CSS y JavaScript. Esto hace que, si visualizamos un archivo HTML desde el navegador que no contiene referencias a ningún otro archivo de otro lenguaje, veamos la información mostrada de forma cruda, sin ningún tipo de estilo ni distribución aplicadas mas allá que los que vienen por defecto en los navegadores.

\subsection{CSS}

Como complemento óptimo del lenguaje HTML, el lenguaje CSS (\textit{Cascading Style Sheets}) define \textbf{la estructura y los estilos utilizados en la presentación} de documentos escritos con HTML.

CSS no es un lenguaje de programación propiamente dicho ni tampoco un lenguaje de marcado, sino que es lo que se ha dado en llamar un \textbf{lenguaje de hojas de estilo}. Su funcionalidad principal es la de definir y mejorar la presentación de páginas web creadas con HTML, creando estilos y formatos que les son de aplicación \cite{mdn:css}.

De la misma forma que HTML, CSS es uno de los lenguajes más usados hoy día en el diseño de los estilos de páginas web, ya sea de forma pura o mediante el uso de distintos frameworks que facilitan su uso.

\subsection{JavaScript}

La funcionalidad principal del lenguaje JavaScript es la de mejorar la experiencia de los usuarios de páginas web, haciendo que las páginas dejen de ser estáticas para pasar a ser \textbf{interactivas y dinámicas}, lo que indudablemente las hace más atractivas e incrementa exponencialmente su funcionalidad \cite{aws:javascript}.

A diferencia de HTML, que es un lenguaje de marcado; y CSS, que es un lenguaje de hojas de estilo; JavaScript es un lenguaje de programación propiamente dicho con \textbf{compilación just-in-time}, es decir, compilación en tiempo de ejecución \cite{mdn:javascript}. Este lenguaje consta de todo lo que podríamos esperar a la hora de escribir código, desde funciones y estructuras de lógica hasta la implementación de programación orientada a objetos.

Al ser un lenguaje de programación con funcionalidad completa, no sólo se usa JavaScript para añadir funciones a las páginas web, sino que hoy en día hay aplicaciones y entornos construidos en JavaScript, como son \href{https://nodejs.org/en}{Node.js}, \href{https://code.visualstudio.com/}{Visual Studio Code}, o \href{https://www.adobe.com/es/acrobat.html}{Adobe Acrobat}. JavaScript también puede ser usado tanto en el \textbf{lado del cliente}, para añadir interactividad a la web; o en el \textbf{lado del servidor}, para manejar la lógica de la aplicación y modificar los datos de esta.

\subsection{Java}

% TODO: Hacer esta parte %

\section{Herramientas de desarrollo}

% TODO: Hacer esta parte %

\subsection{Visual Studio Code}

% TODO: Hacer esta parte %

\subsection{IntelliJ IDEA}

% TODO: Hacer esta parte %

\section{Sistema gestor de bases de datos}

% TODO: Hacer esta parte %

\subsection{MySQL}

% TODO: Hacer esta parte %

\section{Herramientas de gestión}

% TODO: Hacer esta parte %

\subsection{Metodología SCRUM}

% TODO: Hacer esta parte %

\subsection{Git}

% TODO: Hacer esta parte %

\subsection{GitHub}

% TODO: Hacer esta parte %


\capitulo{5}{Aspectos relevantes del desarrollo del proyecto}

En el siguiente apartado se recogerán los aspectos de mayor relevancia que han surgido durante la realización de este proyecto, con el objetivo de comprender mejor el desarrollo y los desafíos encontrados a lo largo del camino.

\section{Inicio del proyecto}

Este proyecto, propuesto por mis tutores, surge como continuidad de otro Trabajo de Final de Grado, el realizado por Mariya Aleksandrova Stroyanova, en el que trabajó para crear una versión web y ampliar las funcionalidades de otro Trabajo de Final de Grado, el realizado por Joseba Fernando Moisén, que consistía en una aplicación de escritorio de la que se podían obtener los valores nutricionales de los alimentos de la base de datos de BEDCA. 

El Trabajo de Mariya, que ha sido el punto de partida, consta de dos aplicaciones web: una encargada de la generación y gestión de empresas, locales, usuarios, menús y platos, que sería la parte que controlarían las empresas encargadas de la gestión de las cafeterías de las distintas facultades de la Universidad de Burgos; y otra aplicación que sería a la que tendrían acceso los clientes, en la cual se pueden realizar informes nutricionales a partir de los menús escogidos por el usuario.

La propuesta para este Trabajo de Final de Grado ha sido la de mejorar estas aplicaciones web, tanto a nivel de infraestructura y servicios, como a nivel de interactividad y usabilidad, procediendo a realizar un rediseño completo de ellas y haciéndolas mucho más amigables a los dispositivos móviles, ya que es el lugar desde donde van a acceder la mayoría de usuarios que visiten la cafetería.

\section{Infraestructura de la aplicación}

Inicialmente, las aplicaciones web parten de un modelo de despliegue de forma manual y en local, ya que requieren de la instalación de varios programas, como MySQL Community, para poder lanzar la base de datos; el SDK de Java, para poder ejecutar el código; así como Eclipse o un IDE similar, para poder lanzar los proyectos de Spring Boot en los que están contenidas las aplicaciones. Esto, como se puede observar en el apartado de Documentación técnica de programación en los anexos de Mariya \cite{tfg-mariya:anexos} es un proceso bastante laborioso, que requiere de la preparación y configuración de un entorno específico para esta tarea, y que no dota a la infraestructura de ningún tipo de escalabilidad ni posibilidad de automatización a la hora de su despliegue.

Además de esto, este modelo de infraestructura es dependiente del sistema operativo y la arquitectura de hardware usadas, lo cual puede suponer un gran inconveniente, como ha terminado siendo el caso, ya que mi equipo personal es un MacBook Pro con procesador Apple Silicon, el cual está basado en una arquitectura ARM. Esto ha supuesto un gran problema a la hora de realizar el despliegue inicial de las aplicaciones, ya que el despliegue está pensado para ser realizado desde un equipo con sistema operativo Windows. 

Primero se intentó solucionar este problema mediante el uso de una máquina virtual, puesto que existen versiones tanto de Windows 10 como de Windows 11 para equipos con arquitectura ARM (cosa no exclusiva de los procesadores de Apple, ya que los propios equipos de Microsoft, los portátiles Surface, actualmente también están dotados de procesadores ARM). Sin embargo, el gran inconveniente surgió a la hora de instalar MySQL, ya que nada más abrir el instalador se obtuvo el siguiente error:

\imagen{MySQL/mysql_installation_error}{MySQL Workbench sólo es compatible con equipos de arquitectura x86$\_$x64}{0.8}

Esto, unido a la motivación de la búsqueda de un despliegue rápido, escalable y moderno, ha terminado en la decisión de realizar un aislamiento completo de los distintos servicios (frontend, backend y base de datos) y contenerizar estos servicios mediante el uso de \textbf{Docker y Docker Compose}, para poder orquestarlos y gestionarlos de forma conjunta. Esto trae numerosas ventajas al proyecto:

\begin{itemize}
  \item \textbf{La infraestructura pasa a ser totalmente independiente del sistema operativo y la arquitectura de procesador}, ya que Docker funciona con todos los sistemas y arquitecturas principales.
  \item \textbf{El despliegue se vuelve mucho más rápido y sencillo}, ya que en prácticamente un minuto y con tan sólo dos comandos puedes estar ejecutando la aplicación con todos sus servicios activos. 
  \item \textbf{Este tipo de infraestructura permite la creación de distintos entornos} (desarrollo, producción, testing) consistentes y sencillos de implementar, ya que permite asegurarse de que la aplicación se va a ejecutar de la misma forma en todos ellos.
  \item \textbf{Cada contenedor se encuentra aislado del resto}, con sus propios recursos y sistemas de archivos, por lo que mejora bastante la seguridad y asegura que un problema en uno de ellos no va a afectar ni canibalizar al resto.
\end{itemize}

Además de esto, puesto que ambas aplicaciones trabajan con el mismo conjunto de datos y no existe realmente un grado de separación necesario como para mantener dos proyectos, también se ha decidido unificar las dos aplicaciones web en una, debido a que esto facilita bastante su desarrollo, elimina código repetido y mantiene mayor homogeneidad.

\section{Lógica de negocio}

En el proyecto original, ambas aplicaciones siguen el patrón \textbf{Modelo-Vista-Controlador (MVC)} en todo su desarrollo, tanto para el manejo de datos como la interfaz, ya que todo forma parte de un mismo monolito. Sin embargo, como en este proyecto se ha decidido hacer una separación total de los servicios, el enfoque adoptado ha sido el de \textbf{API-Driven Development}, en el que primero se implementa esta lógica de negocio siguiendo el patrón MVC, y en la capa del Controlador se implementa una API REST que proporciona los endpoints necesarios para que el frontend interactúe con esta. 

Esto ha permitido abstraer a la aplicación web de tener que interactuar directamente con la base de datos, lo que ahorra el uso de queries SQL complejas y permite prescindir de APIs como JDBC, ya que se va a trabajar directamente con objetos JSON mucho más amigables de manejar y representar en la capa frontend.

Para la implementación de este servicio se ha mantenido el uso de \textbf{Spring Boot}, al igual que en el proyecto original, lo que ha permitido tener el modelo original como referencia a la hora de realizar esta reimplementación.

\subsection{Modelo de datos}

Puesto que tanto el proyecto de Mariya como el actual siguen el modelo MVC, las capas del Modelo de datos se asemejan, ya que en ambos se mantiene un modelo puramente relacional. Sin embargo, esta transformación en la capa del Controlador, así como los cambios realizados a la hora de obtener la información de los alimentos, han hecho que haya varias diferencias entre los dos:

\imagen{DB/mariya_data_schema}{Modelo de base de datos del TFG del que se parte \cite{tfg-mariya:anexos}}{0.7}

Como los datos van a ser accedidos y manipulados directamente desde la API, se ha eliminado la necesidad de tener vistas en la base de datos, ya que los propios endpoints son los encargados de mostrar sólo la información que nos interese, permitiendo además mayor personalización y granularidad en los resultados obtenidos.

Además de esto, como en la sección de Alimentos se ha pasado de un modelo en el que los datos sólo se podían almacenar de forma local, usando los datos de BEDCA, a un modelo híbrido que permite tanto el almacenamiento local como la interacción con la API de Nutritionix, esta parte del modelo de datos se ha rediseñado completamente, intentando hacerla más simple y fácil de entender.

\imagen{DB/nutrimenu_data_schema}{Modelo de base de datos actual}{0.7}

Para la implementación de la base de datos se ha mantenido \textbf{MySQL}, ya que al seguir manteniendo un modelo de bases de datos relacional no he visto necesidad de modificarlo.

\section{Frontend}

Este es el servicio que más cambios ha sufrido, puesto que prácticamente ha sido recreado desde 0. Para el desarrollo de esta aplicación web se ha usado la librería \textbf{React}, debido a que actualmente es de las opciones más populares en el desarrollo frontend, y su modelo basado en Componentes me pareció muy buena solución para el desarrollo de este proyecto, puesto que prima bastante la reutilización y la modularidad. 

Sin embargo, como React no dispone de forma nativa de ningún sistema de enrutamiento, puesto que su modelo principal es el de Single Page Applications (aplicaciones que reescriben el contenido mostrado de forma dinámica, en vez de recargar la página completa), el proyecto da bastante uso a otra librería pensada para ser usada junto a React, que es \textbf{React Router}. Esta librería ha permitido gestionar las distintas partes de la aplicación sin ningún problema, pues al ser una aplicación con distintas secciones bien diferenciadas, siguiendo el modelo de las SPA sería mucho más complicado de separar y limitaría bastante las posibilidades.

Además de esto, \textbf{React Router} incluye otras funciones bastante aprovechadas en este proyecto, como los \textbf{Outlets}, que permiten crear un componente plantilla y que los componentes hijos se rendericen dentro de ese componente padre (usados en varias partes, como por ejemplo en los layouts de la barra de navegación y los paneles de gestión de los elementos, como se puede ver en la figura \ref{fig:App/outlets_demo}); o los \textbf{loaders}, que son funciones que permiten realizar llamadas asíncronas a una API antes de que se renderice el componente que va a usar esos datos (usados cada vez que se muestra información sobre un alimento, plato o menú), lo cual acelera bastante la carga.

\imagen{App/outlets_demo}{Aplicación renderizando sólo el contenido de los componentes padre}{1}

El desarrollo se ha realizado mediante el uso de \textbf{JavaScript}, lo que ha supuesto posiblemente el mayor reto de todo el proyecto, pues previo a la realización de este trabajo no había tenido ninguna experiencia con este lenguaje.

Como he usado React, a la hora de escribir el código referente a la interfaz he usado \textbf{JSX}, una extensión de JavaScript que permite integrar tanto HTML como JavaScript como si fueran un sólo lenguaje.

Para la construcción y compilación del frontend inicialmente se comenzó usando \textbf{Create React App}, ya que hasta hace poco era la opción recomendada por los propios desarrolladores de React, pero puesto que es un proyecto que ha dejado de ser mantenido, y los paquetes que contiene están quedando desactualizados, finalmente se ha cambiado por \textbf{Vite}, una de las herramientas de compilación de JavaScript más eficientes de la actualidad. Esta herramienta ofrece distintas optimizaciones a la hora de ejecutar y compilar nuestro código que hacen que la aplicación web sea lo más rápida posible, lo que se nota bastante al usarla, ya que todo es prácticamente instantáneo.

En el estilado de la aplicación se ha trabajado principalmente con \textbf{Bootstrap}, ya que anteriormente ya había realizado algún proyecto dando uso de este framework y tenía algo de experiencia, y ha sido de bastante ayuda al crear una interfaz responsive y visualmente agradable, adaptada tanto a pantallas de dispositivos móviles como a pantallas de ordenador. 

Existe una librería llamada \textbf{React Bootstrap} que integra todos los elementos de Bootstrap dentro de React, para trabajar con ellos como si fueran componentes de React, por lo que me he ahorrado el uso de bastantes clases de CSS.

Por último, para la representación de las tablas y los gráficos en las vistas de información nutricional (ejemplo en la figura \ref{fig:App/charts_and_tables}), he usado las librerías \textbf{AG React Data Grid} y \textbf{AG React Charts} respectivamente, debido a que requieren de muy poco código para conseguir los resultados buscados, y me permiten actualizar los datos de forma dinámica, algo que necesitaba para poder cambiar los datos nada más se seleccione un plato distinto en la vista del informe nutricional de un menú.

\imagen{App/charts_and_tables}{Distintas tablas y gráficos usados en la aplicación}{1}

\section{Fuente de datos nutricionales}

Las aplicaciones originales hacían uso de la base de datos de BEDCA, la Base de Datos Española de Composición de Alimentos, para obtener toda la información relacionada con los alimentos y sus componentes nutricionales. 

El gran inconveniente de esto es que esta base de datos no dispone de un punto de acceso al que consultar los datos mediante una API o similar, así que era necesario mantener la base de datos en local. Esto, además de que impide obtener la última información nutricional disponible, puede suponer un costo en almacenamiento y rendimiento importante, ya que todos los alimentos y su información deben estar guardados en el servidor obligatoriamente.

Para solucionar esto, se ha decidido dar uso de la \textbf{API de Nutritionix}, que dispone de una base de datos específica para alimentos en castellano. Esto permite almacenar sólo los alimentos que vayan a ser usados, mientras que el resto quedan disponibles con tan sólo realizar una búsqueda en la aplicación (como se ve en la figura \ref{fig:App/nutritionix_search}). Además, para suplir aquellos alimentos que pueden no estar presentes en esta nueva base de datos, también se ha dado la opción de crear los alimentos a mano, para que en caso de que sea necesario un alimento esto no sea un factor limitante. 

\imagen{App/nutritionix_search}{Búsqueda de un alimento mediante el uso de la API de Nutritionix}{0.8}

El mayor inconveniente que he encontrado con el uso de esta API es el hecho de que no contiene información nutricional sobre alérgenos, quedando esto como objeto de mejora del proyecto.

\section{Automatizaciones}

Como la infraestructura de la aplicación dispone de varios servicios independientes, es posible que en algún momento realice un cambio en uno de ellos y sin darme cuenta esto afecte a otro distinto, haciendo que algún elemento no funcione como debería. Es por ello que he realizado una implementación de integración continua mediante el uso de \textbf{GitHub Actions}, aprovechando que el despliegue se realiza con Docker y Docker Compose. 

De lo que se encarga esta automatización es que, cada vez que se hace un push de algún commit al repositorio de GitHub, se ejecutan una serie de pasos (como podemos ver en la figura \ref{fig:App/github_actions_execution}) para desplegar los contenedores de Docker en una máquina nueva, y le da un minuto desde que todos los contenedores están corriendo para que cada servicio devuelva un \textit{healthcheck} indicando que todo está funcionando correctamente. En caso de que alguno de los servicios no devuelva este mensaje, el propio GitHub se encarga de enviar un mensaje indicando que algo ha fallado, además de indicarlo en el propio commit.

\imagen{App/github_actions_execution}{Listado de pasos ejecutados por la acción de GitHub Actions}{0.8}
\capitulo{6}{Trabajos relacionados}

\begin{comment}
	
Este apartado sería parecido a un estado del arte de una tesis o tesina. En un trabajo final grado no parece obligada su presencia, aunque se puede dejar a juicio del tutor el incluir un pequeño resumen comentado de los trabajos y proyectos ya realizados en el campo del proyecto en curso. 

\end{comment}

\capitulo{7}{Conclusiones y Líneas de trabajo futuras}

Todo proyecto debe incluir las conclusiones que se derivan de su desarrollo. Éstas pueden ser de diferente índole, dependiendo de la tipología del proyecto, pero normalmente van a estar presentes un conjunto de conclusiones relacionadas con los resultados del proyecto y un conjunto de conclusiones técnicas. 
Además, resulta muy útil realizar un informe crítico indicando cómo se puede mejorar el proyecto, o cómo se puede continuar trabajando en la línea del proyecto realizado. 



\bibliographystyle{plain}
\bibliography{bibliografia}

\end{document}
